\scanpage{283}
\mspage{199}
% archive sheet number: 282
\deleted{ce foncteur par $C_{/\unsure{x|X}}$,}

\deleted{\diagram{\begin{tikzcd}[ampersand replacement=\&,column sep=large]
I\arrow[r] \& C_{/X}\arrow[r,dashed] \& C\arrow[r,dashed] \& M{,}
\end{tikzcd}}}

\deleted{et on a vu (\unsure{cf} 4.4.2 p.~\unsure{74|79|71}) que ce foncteur est cofinal.}

En effet, le foncteur canonique $M\longrightarrow C^{\wedge}$ est pleinement fidèle, et on applique le critère \unsure{c)} de 4.8.2.

Pour achever de prouver 4.8.1, il reste donc à établir l'implication

\[ \mathrm a)\Longrightarrow c'. \]

\deleted{Soit donc $S\subset M$}

Supposons donc que $M$ satisfasse à $L_{\unsure{\pi|\kappa|k}_0}$, et admette une petite sous-catégorie str. génératrice $S$, formée d'objets $\pi_0$-accessibles. \deleted{L'hyp. $L_{\pi_0}$ \unsure{signifie|dit} que} \inserted{On} \unsure{trouve|tient} un foncteur canonique

\[ \mathrm{Ind}_{\pi_0}S\longrightarrow M, \]

prolongeant l'inclusion \deleted{canonique} $S\hookrightarrow M$, de façon à commuter aux $\varinjlim$ \unsure{grdes} devant $\pi_0$, et du fait que $S\subset M_{\pi_0}$, ce foncteur est \uline{pleinement fidèle}. Mais comment prouver qu'il est \unsure{ess. surjectif} ? \unsure{Ca|Ce} ne sera d'ailleurs pas vrai \deleted{\unsure{à bien}} sans autre hypothèse sur $S$. \deleted{\unsure{en fait}} L'idée naturelle est de saturer $S$ pour les $\varinjlim$
